\section{Decode the format's byte order}
The format puts the most significant byte of each 16-bit field first. If the two
unsigned bytes are \(h\) and \(l\), the decoded integer is
\[
 u=256h+l=(h\ll8)\mathbin{\mathrm{OR}}l.
\]
For sequence bytes \texttt{12 34}, this gives hexadecimal \texttt{1234}, or
\(18\cdot256+52=4660\). For \texttt{80 ff}, it gives 33023. The latter remains
positive: an unsigned 16-bit field ranges from zero through 65535.

\texttt{lbu} loads one byte and extends it with zeros, yielding zero through 255.
The shifted high byte occupies bits 8 through 15; the low byte occupies bits zero
through seven, so OR combines them without overlap. Both values fit a positive
signed word. Use word arithmetic for these bounded field values and full-width
arithmetic for addresses.

The default MARS memory configuration is little-endian. A halfword load would
interpret the low-address byte as the less significant one and would also impose
alignment requirements. The input pointer can be unaligned. Separate byte loads
make the specified order explicit and work at every allowed pointer offset. The
four output words use ordinary native word storage. Their values are decoded
integers; the output record is not another copy of the wire-format header.

Do not confuse a signed source literal with a signed field interpretation. The
starter writes payload bytes 128 and 255 as \texttt{.byte -128, -1}, which preserves
their eight-bit patterns without an out-of-range signed-byte warning. Loading those
patterns with \texttt{lbu} yields 128 and 255. A debugger's signed decimal view may
show different numbers for the same stored bits.
